\section{Manuscript audit and proposed corrections}
\label{sec:audit}

The audit is tied to the fixed revision in Section~\ref{sec:scope}.
Some objections recorded in earlier research packages have already been
corrected in that revision. They should not be reported again as newly
discovered defects. The following changes are supported by the proofs
and direct source inspection in this article.

\subsection{A confirmed even-weight normalization error}

In \src{chapters/03-algebraic.tex}, equation
\texttt{bloch:eq:Pm} defines
\[
 P_m(x)=\operatorname{Re}_m\left(
  \sum_{j=0}^{m-1}\frac{2^jB_j}{j!}
       (\log|x|)^j\Li_{m-j}(x)\right),
 \qquad
 \operatorname{Re}_m=
 \begin{cases}\operatorname{Re},&m\text{ odd},\\
 \operatorname{Im},&m\text{ even}.
 \end{cases}
\]
For $0<x<1$, every term inside the parentheses is real. Consequently
\begin{equation}\label{eq:even-P-zero}
 P_{2j}(x)=0\qquad(0<x<1).
\end{equation}
This is also an immediate consequence of the conjugation law printed
in the same chapter. The paragraph immediately following
\texttt{bloch:eq:goldenP3} and \texttt{bloch:eq:sgP3} nevertheless says
that the higher golden ladders become the stated zeta multiples under
the same modification, and explicitly includes the weight-six multiple.
That assertion is incompatible with \eqref{eq:even-P-zero}.

The correction is local. Keep the two displayed weight-three formulas.
Explain that their lower-weight logarithmic terms cancel under $P_3$.
At other odd weights, carry out the full modified-polylogarithm
calculation and state the proof status of the resulting formula. At
even weight, keep the ordinary $\Li_{2j}$ identity in its ordinary
normalization; it cannot become a nonzero real zeta multiple of this
$P_{2j}$ at real arguments. A different real-valued normalization would
need to be explicitly defined and checked. The file
\src{integration/proposed_P_even_correction.tex} contains replacement
prose ready to paste into the indicated location. The associated patch
is restricted to this confirmed error.

This correction does not disprove the numerical ordinary-polylogarithm
formulas at weights six and eight. It corrects their asserted
interpretation. The distinction matters because a vanishing
single-valued evaluation cannot retain the nonzero zeta value that the
paragraph assigns to it.

\subsection{Three numerical statements can now be promoted}

The equation \texttt{gauss:eq:S4-closed} in
\src{chapters/04-depth.tex} is an exact theorem by
Theorem~\ref{thm:s4}. The surrounding paragraph currently places the
whole displayed group in a numerically supported category unless an
analytic derivation is supplied. Add a specific reference to the present
proof and certificate for this equation. This does not automatically
promote neighboring identities whose proofs are not given here.

In the subsection on the reciprocal minimal Pisot number in
\src{chapters/03-algebraic.tex}, the displayed plastic dilogarithm and
trilogarithm formulas now have the explicit proofs supplied in the
preceding sections. Their normalization agrees exactly with the
manuscript: $6/4=3/2$, $3/9=1/3$, and $6\lambda^3/3!=\lambda^3$.
The adjacent plastic tetralogarithm, which introduces the exponents
$7$ and $14$, is a separate identity. No proof of that formula follows
merely by citing the two lower-weight results.

Several higher golden formulas are described using language such as
``hold'' or ``correct'' immediately after residual calculations.
Where no independently traceable analytic proof is supplied, prefer
``numerically verified to the stated precision.'' A residual, even an
extremely small one reproduced by different evaluators, establishes
neither an exact equality nor its uniqueness. This is a proof-status
clarification, not an assertion that those formulas are false.

\subsection{What the classification does and does not imply}

The pure-complement classification is exhaustive. It permits the exact
statement that the reciprocal plastic family is the only real positive
cyclic-power family with two different complement pairs. It does not
bound the maximal weight of a polylogarithm ladder. A relation such as
\[
 \prod_j(1-r^j)^{n_j}=r^N,
 \qquad n_j,N\in\mathbb Z,
\]
may use many factors that are not themselves pure powers of $r$.
Nor does the classification exclude additional units or nonmonomial
arguments in the field. These possibilities are central to higher
ladders and must remain distinct from the theorem proved here.

The degree bound is also specific: it makes the search for
$r^a+r^b=1$ complete at fixed degree. A finite coefficient-height box
for arbitrary algebraic arguments has no comparable completeness
consequence without an additional height theorem.

\subsection{Analytic qualifications needed for integration}

There are three limits on the gamma and harmonic statements worth
preserving verbatim when they are integrated.

First, the late reflected-coefficient series is divergent for every
fixed $n$. Its finite truncations form an asymptotic expansion, and
the growing-order theorem supplies a rigorously controlled truncation.
The displayed bound is unsigned; eventual alternation of coefficients
does not establish the sign of the optimized remainder.

Second, $m$ is fixed in the reflected-moment theorems. Reflection
symmetry $M_{n,m}=M_{m,n}$ does not turn these statements into a theorem
when both indices grow proportionally. In the harmonic theorem,
$\alpha=p/h$ stays in a compact subset of $(0,\infty)$. Substitution
of $\alpha=\log h+c$ into its error term is not justified by that
compact-uniform assertion alone.

Third, the Appell identity concerns canonical formal coefficients
defined by the generating function. It gives an unconditional formal
version of the inspected author preprint's Conjecture~5.5. It does
not prove transcendence or uniqueness of arbitrary numerical
expressions involving logarithms and zeta values. The attribution is
to the explicitly inspected 2010 preprint, without asserting that the
formulation remained open in every later version of the literature.

\subsection{Suggested integration sequence}

The new material can be integrated in four independently reviewable
units. Put the complementary-power classification and the two plastic
proofs beside the existing algebraic ladders. Put the $S_4$ theorem,
its word convention, and its certificate beside the Gaussian weight-five
identities. Add the proportional-depth theorem to the alternating-harmonic
report and cross-reference it from the depth chapter. Add the reflected
moment results to the gamma-moment material, retaining the earlier
unreflected theorem as the special case $m=0$.

The article and all certificates are self-contained within this package.
The supplied narrow patch uses the audited source context and can be
reviewed before application. No repository branch, commit, or remote
publication was created as part of this delivery.
